
\subsubsection{4.1 Experimental Questions}

\textbf{Q1 (h-e1).} Does ratio reward produce detectably different GRPO advantage variance than binary reward in groups representative of early-training conditions?

\textbf{Q2 (h-m1).} Does the gradient signal difference (Q1) translate to measurable policy-level performance improvements after 1,000 GRPO training steps?

\subsubsection{4.2 Metrics}

\textbf{Advantage variance.} Var({r\_i $-$ mean(r)}) computed per group, measuring the gradient signal available from that group.

\textbf{FractionPartialCallback.} Fraction of completions per batch with 0 < k < n test cases passing. When fraction\_partial = 0, ratio and binary rewards assign identical values to all completions.

\textbf{clipped\_ratio.} Fraction of completions reaching max\_new\_tokens limit. clipped\_ratio = 1.0 indicates all completions are truncated at the token limit, precluding test execution.

\textbf{HumanEval pass@1.} Greedy-decode pass@1 on 164 problems \citep{Chen2021HumanEval}, evaluated at training checkpoints.

\subsubsection{4.3 h-e1 Components}

\begin{enumerate}
\item \textbf{Synthetic group computation.} Direct computation of advantages for group [0,0,1,2,0,3,0,0] (8 completions, 5 test cases) under both reward functions.
\item \textbf{Scale simulation.} 1,000 GRPO groups (G = 8, n = 5, p\_pass = 0.1 per test); record fraction where ratio $\neq$ binary signal.
\item \textbf{Smoke test.} 3 actual GRPO training steps; verify exit code 0, gradient norm logging, checkpoint saving.
\item \textbf{Background training run.} 151 steps launched (PID 2605366); gradient norm CI computed via bootstrap.
\end{enumerate}

\subsubsection{4.4 h-m1 Design}

Full GRPO training for 1,000 steps under both conditions simultaneously on separate GPUs. Checkpoints at steps 200, 400, 600, 800, 1,000. HumanEval evaluated at all checkpoints. FractionPartialCallback fires at every step as the primary diagnostic for ratio reward degeneracy.

The experiment ran for approximately 208 logged training steps before being terminated and routed to Phase 0 redesign, based on the evidence that both conditions had identical zero-reward signals throughout.

